\begin{claim}{C.53}{Proved}{Global quotients and finitely many operator representatives}
Use the ambient-null based plane loop $f_t$, its generated group
$G$, the commuting operator group $H$, and $L=\overline G$ of
C.42. For every open $G$-orbit $O$, choose $x\in O$ and
put $J=L_x$. Define $\mathscr F_O^{\rm all}$ to be the family
of fiber partitions on the common universal covering of $O$
of all real lifts of
\[
 \Phi_z^x:O\longrightarrow Gz,\qquad \Phi_z^x(lx)=lz,
 \quad z\in\operatorname{Fix}(J),\quad
                  Gz\text{ nonopen and nontrivial}.
                                                        \tag{C.53.1}
\]
The target points need not belong to $\partial O$.
Let $U^\dagger$ be the union of all open orbits with either
$\dim Q_O=2$ or $|\mathscr F_O^{\rm all}|\geq3$, and set
$X^\dagger=\overline{U^\dagger}$.
There is a jointly continuous based contraction on
$X^\dagger$, with jointly continuous inverse and exact
$H$-commutation, extending the prescribed contraction on
$X_*$ in C.50. On every nonopen orbit it has the
corresponding exponential value of C.52.1.

Now choose finitely many open orbits
$O_1,\ldots,O_m$ disjoint from $X^\dagger$, each with compact
$O_j/Q_{O_j}$, and put
\[
 U=\bigcup_{j=1}^m\ \bigcup_{h\in H}hO_j,
 \qquad X=\overline U.
                                                        \tag{C.53.2}
\]
The original loop has a based contraction through
homeomorphisms of $X^\dagger\cup X$, jointly continuous
with its inverse, commuting with $H$ and equal to the
specified exponential contraction on $X^\dagger$.
On every nonopen orbit in $X$ its value is exactly the
hybrid prescription $R^B_{s,t}$ of C.52.1.
This assertion allows infinitely many different open
components in (C.53.2), with arbitrary accumulation.
It retains the stated finite number of $H$-representatives
and compactness of their actual operator quotients.
\end{claim}

\begin{proof}
\emph{The global quotient family and its exact comparisons.}
C.42 is already stated for every $z\in\operatorname{Fix}(J)$
with nontrivial line or circle orbit; it imposes no
frontier condition on $z$. It gives the actual continuous
open map (C.53.1), its intrinsic target topology, the
common covering construction, and its proper fiber lines.
C.43--C.45 use these maps and their transitive compatible
action to give the connected real fibers, product
coordinates and the exact three pairwise coordinate
domains. Thus each of those conclusions applies to
(C.53.1) with the same maps.

Changing the source basepoint also preserves the complete
family. If $y=ax$ for $a\in L$, then
$L_y=aJa^{-1}$ and, for $z\in\operatorname{Fix}(L_y)$,
\[
                \Phi_z^y(ly)=lz
                     =\Phi_{a^{-1}z}^x(lax).
                                                        \tag{C.53.3}
\]
The targets are the same actual $G$-orbit, because
$Lz=Gz$ for each such quotient, as proved in C.42.
This proves the asserted identification of partitions,
including the original target coverings.

For $h\in H$ stabilizing $O$, choose $a_h\in L$ with
$a_hx=hx$. Exact commutation gives
$L_{hx}=L_x$, hence $a_hJa_h^{-1}=J$, and
\[
                   \Phi_z^x h=\Phi_{a_hz}^x.             \tag{C.53.4}
\]
The point $a_hz$ is fixed by $J$ and lies on the same
nonopen orbit as $z$. Therefore (C.53.4) permutes the
entire family $\mathscr F_O^{\rm all}$. For $hO=O'$, use
basepoint $hx$ in $O'$; then
\[
               \Phi_{hz}^{hx}h=h\Phi_z^x.              \tag{C.53.5}
\]
These equalities identify the families under operator
transport and prove that their cardinalities are
well defined and $H$-invariant.

\emph{Three global quotients give the same actual Lie group.}
Choose three distinct partitions in
$\mathscr F_O^{\rm all}$. The full three-foliation group
of C.47 is constructed exclusively from the real maps
in C.43--C.45 and their common transitive action.
Its proper-pair or ordered-loop argument therefore
applies to these exact maps and gives the same locally
compact Lie group of dimension two or three.
The based target lifts in C.49.2 are loops: for circle
targets C.29 uses the ambient-null plane loop and two
points of that circle, independently of the position
of $O$; real target endpoints are the identity.
The Baire transitivity proof C.49.3, its isotropy
calculation C.49.4--C.49.7 and its closedness argument
now apply to the same based source lifts. Their actual
generated group $N$ is consequently one of
$\mathbb R^2$, $\operatorname{Aff}^+(\mathbb R)$, or
$\mathcal A_2$, with exactly the topology, global
exponential and unique logarithm proved in C.49.

We verify the action comparison for every global
quotient, rather than only for frontier targets.
Use the closed Polish group
\[
 E_O=\{(l,n)\in L\times N:pn=(l|_O)p\}.
                                                        \tag{C.53.6}
\]
Its projection to $N$ is onto, since every $n$ is an
actual word in the lifted loop; the identical global
word supplies $l\in G$. It is an open homomorphism
of Polish groups. If $(l,1)$ is in its kernel, $l$
fixes $O$ pointwise. Equivariance of (C.53.1) then
makes $l$ fix its entire image $B=Gz$ pointwise:
\[
                         l\Phi_z^x(u)=\Phi_z^x(lu)
                                      =\Phi_z^x(u).
                                                        \tag{C.53.7}
\]
C.42 proves continuity of $L\to\operatorname{Homeo}(B)$
in the intrinsic compact-open topology. Its restriction
to $E_O$ therefore factors continuously through the
open quotient to an actual action of $N$ on $B$.
That image is exactly $G|_B$: actual loop words prove
both inclusions. It is metabelian and hence is the
translation, affine, or rotation Lie image of C.37.
The resulting continuous homomorphism and its unique
based covering lift give exactly
\[
 \Phi_z^x K^O_{s,t}=R^B_{s,t}\Phi_z^x,
 \qquad K^O_{s,t}=
        \exp_N\bigl((1-s)\log_N\widetilde f_t\bigr)
                       \text{ descended to }O.         \tag{C.53.8}
\]
Here $R^B$ is the exponential branch of C.52.1,
with the full based real angle for a circle.
The original source covering, every original quotient,
and every restriction kernel are retained.

\emph{The enlarged closed exponential region.}
For a dimension-two proper-operator orbit use the
based lifted contraction of C.36--C.37. For an orbit
with three global partitions use (C.53.8).
Where both constructions apply, their source groups
are the same actual group generated by the same
based lifted loop, as in C.49, so the maps agree.
Operator conjugation transports these groups and
intertwines their exponential maps, by (C.53.5) and
the based covering identity C.49.16.

Every source group just specified has the uniform
fractional-word bound $1913$: C.37 gives the stronger
bound three for the first two groups, and C.50.1--C.50.14
give the stated bound for $\mathcal A_2$.
Each word is a word in the original based path and
its inverse. In particular the full argument
C.50.15--C.50.19 applies to this enlarged union with
precisely the same compact set of actual global word
isotopies. We spell out the transfer of its remaining
inputs. All the source images are metabelian, so
each global commutator-of-commutators fixes
$U^\dagger$, hence its closure. At a nonopen limit
orbit the actual transitive image is therefore the
same translation, affine or rotation group classified
in C.37. For every integer $n$, the residual after
removing $\lfloor n(1-s)\rfloor$ endpoint powers is a
fractional power in the specified simply connected
source group, and thus has the same uniform word
bound. On a line the bounded residual logarithms
force the exponential limit. On a circle the source
residual word loop is null in that same simply
connected group; its exact two-point winding passes
to the limiting orbit, giving the real equality
C.50.19, including its integer part. These are all
inputs of C.50's limit identification and joint
continuity proof. Applying it also to the inverse
loop gives a jointly continuous inverse on
$X^\dagger$.

The original $U_2\cup U_3$ lies in $U^\dagger$.
On that original union the maps are the same
source fractional powers. Continuity extends this
equality to $X_*$. The operator identities also
extend from the dense union. This proves the first
assertion and fixes its prescription for the rest
of the proof.

\emph{Compatible intrinsic maps on the remaining representatives.}
Each $O_j$ in (C.53.2) has
$|\mathscr F_{O_j}^{\rm all}|\leq2$ and
$\dim Q_{O_j}\leq1$, by its disjointness from
$X^\dagger$. For a family of size one or two,
apply the construction of C.52 to this full global
family. Here is the exact transfer. Its source
coordinates and operator normal forms use
C.42--C.45, already verified above for (C.53.1).
Its full-partition invariance is (C.53.4).
The base-family replacement C.52.5 is natural under
every target reparametrization and exchange, and
retains every affine operator equation C.52.6.
No step of those intrinsic constructions uses the
position of the target point on a frontier.
They therefore give a based intrinsic contraction
$K^j_{s,t}$, jointly continuous with its inverse,
commuting with $H_{O_j}$ and with its actual
restricted closure $Q_{O_j}$, and satisfying
\[
                    \Phi_z K^j_{s,t}
                             =R^{Gz}_{s,t}\Phi_z
                 \quad\text{for every map in (C.53.1)}.
                                                        \tag{C.53.9}
\]
For an empty family, C.32--C.33 give the based
intrinsic contraction in that exact proper-operator
centralizer. Every constant equivariant quotient
also satisfies (C.53.9), since its target is fixed.

Discard repetitions among the finitely many
$H$-orbits represented by $O_j$, and transport
$K^j$ to each $hO_j$ by $hK^jh^{-1}$.
Commutation with the full stabilizer makes this
well defined. Its union gives a jointly continuous
homeomorphism family $K$ on the open set $U$;
components are open, and each carries its intrinsic
family and inverse. It commutes with $H$ exactly.

\emph{One compact comparison set for every accumulating sequence.}
For each representative, compactness of
$O_j/Q_{O_j}$ supplies a compact $C_j\subset O_j$
with $Q_{O_j}C_j=O_j$, by C.51's finite-cover
argument. Suppose
\[
 x_n\in U,\quad x_n\longrightarrow z\in X\setminus U,
                  \qquad(s_n,t_n)\longrightarrow(s,t).
                                                        \tag{C.53.10}
\]
The orbit $Gz$ is nonopen. Indeed an open orbit
containing $z$ would contain every sufficiently
late $x_n$, and so would itself be one of the
components in $U$, contradicting $z\notin U$.
There are only finitely many representatives.
On a subsequence, write
\[
                  x_n=h_nq_ny_n,
 \quad h_n\in H,\quad q_n\in Q_{O_j},\quad
                  y_n\in C_j,\quad y_n\longrightarrow y\in C_j
                                                        \tag{C.53.11}
\]
for one fixed $j$. Fix the compact subset of $O_j$
\[
 S_j=C_j\ \cup\!
       \bigcup_{(u,v)\in[0,1]^2}
       \bigl(K^j_{u,v}(C_j)\cup(K^j_{u,v})^{-1}(C_j)\bigr).
                                                        \tag{C.53.12}
\]
Both joint continuity assertions prove compactness.

Choose actual $a_n\in H_{O_j}$ so that
\[
 \sup_{w\in S_j}
          |h_na_nw-h_nq_nw|<1/n,
 \qquad g_n=h_na_n\in H.
                                                        \tag{C.53.13}
\]
For each fixed $n$, the first condition is an
open neighborhood of $q_n$ in
$\operatorname{Homeo}(O_j)$: composition with the
fixed ambient homeomorphism $h_n$ and restriction
to the compact set $S_j$ are continuous.
The defining density of the actual restrictions
$H_{O_j}\to Q_{O_j}$ supplies $a_n$.
Thus the estimate is made \emph{after} the possibly
distorting operator $h_n$; no uniform modulus of
continuity for the sequence $h_n$ is assumed.

The open $L$-orbit map supplies $l_n\to1$ with
$l_ny=y_n$, exactly as in C.51. Since $g_n$
commutes with $L$, (C.53.10)--(C.53.13) give
\[
                       g_ny=l_n^{-1}g_ny_n\longrightarrow z.
                                                        \tag{C.53.14}
\]
Every element of $L_y$ fixes this limit.
The evaluation homeomorphism C.42.1 therefore gives
\[
                g_n|_{O_j}\longrightarrow\Phi_z^y
                     \quad\text{uniformly on compact sets}.
                                                        \tag{C.53.15}
\]
The exact basepoint change (C.53.3) identifies this
map with one of the global quotients in (C.53.1)
when its target is nontrivial; otherwise it is
constant. This is the reason the enlarged family
is required in this argument.

Transport and commutation with $Q_{O_j}$ now yield
\[
 \begin{split}
 K_{s_n,t_n}(x_n)&=h_nq_nK^j_{s_n,t_n}(y_n),\\
 \bigl|K_{s_n,t_n}(x_n)
              -g_nK^j_{s_n,t_n}(y_n)\bigr|&<1/n.
 \end{split}                                            \tag{C.53.16}
\]
The inner points converge in $S_j$ to $K^j_{s,t}(y)$.
Equations (C.53.15) and (C.53.9) consequently give
\[
                   K_{s_n,t_n}(x_n)\longrightarrow
                         \Phi_z^yK^j_{s,t}(y)
                           =R^{Gz}_{s,t}(z).            \tag{C.53.17}
\]
Every subsequence admits a further subsequence of
this form, by the finite number of $j$ and the
compactness of their $C_j$. The value in (C.53.17)
is independent of every such choice. Hence the
entire sequence has that limit. If a representative
has an empty global family, (C.53.15) also proves
that every limit reached through its translates
is a fixed point.

Define the extension on $X\setminus U$ to be
$R^{Gz}_{s,t}(z)$. The diagonal approximation
C.52.9 now applies with the dense set $U$ in $X$:
for each varying boundary point and fixed parameter
pair choose a point of $U$ close to it and with
close image, using (C.53.17). Applying (C.53.17)
to the resulting sequence proves joint continuity
on all of $X$. This argument does not require
prior transverse continuity of the hybrid rule.

Use the inverse images already present in $S_j$
and the exact inverse version of (C.53.9) to
repeat (C.53.13)--(C.53.17) for inverse maps.
The same diagonal approximation then proves
joint inverse continuity on $X$.
The maps preserve $U$ and each nonopen limit
orbit, with the displayed maps as mutual
inverses. Thus they are homeomorphisms of $X$.
Their based-edge identities and $H$-commutation
extend from $U$ by continuity; the restriction
to every nonopen orbit is the prescribed $R^B$.

Finally, $U$ is disjoint from $X^\dagger$.
Every point of $X\cap X^\dagger$ lies on a
nonopen orbit; otherwise its open orbit would
meet $U$ and be a component of $U$.
The metabelian identity on $X^\dagger$ and the
first part of the proof show that both
contractions there use precisely the same
exponential branch of C.52.1, with the same
based covering path. They and their inverses
agree on the full intersection. Pasting the two
closed sets $X$ and $X^\dagger$ gives the
asserted contraction on their union, preserving
the prescribed family on $X^\dagger$.

The finite-representative argument does not
supply compact comparison sets for arbitrary
noncocompact quotients or infinitely many
$H$-inequivalent representatives. Nor does it
extend the family through an exterior region
of nonopen orbits. These remain exact parts
of the general C.12 obligation.
\end{proof}
